\documentclass[aspectratio=169]{beamer}  % 16:9 aspect ratio

% Use a clean theme as base
\usetheme{default}
\usecolortheme{default}
\useoutertheme{miniframes}
\useinnertheme{circles}
\usepackage{booktabs}
\usepackage[backend=bibtex,sorting=none]{biblatex}
\addbibresource{ref.bib} 
\setbeamerfont{footnote}{size=\tiny}

% Custom colors from HKUST logo
\definecolor{hkustblue}{RGB}{0, 51, 119}    % Navy blue from logo
\definecolor{hkustgold}{RGB}{180, 141, 61}  % Golden brown from logo
\definecolor{lightgray}{RGB}{236, 240, 241}


% Customize the appearance
\setbeamercolor{structure}{fg=hkustblue}
\setbeamercolor{background canvas}{bg=white}
\setbeamercolor{normal text}{fg=hkustblue}
\setbeamercolor{frametitle}{fg=hkustblue,bg=white}
\setbeamercolor{itemize item}{fg=hkustgold}
\setbeamercolor{itemize subitem}{fg=hkustgold}
\setbeamercolor{block title}{fg=white,bg=hkustblue}
\setbeamercolor{block body}{fg=hkustblue,bg=lightgray}
\setbeamercolor{title}{fg=hkustblue}
\setbeamercolor{author}{fg=hkustgold}


\setbeamertemplate{title page}
{
  \begin{centering}
    \vspace{1cm}
    {\usebeamercolor[fg]{title}\LARGE\inserttitle\par}
    \vspace{0.8cm}
    {\usebeamercolor[fg]{author}\large\insertauthor\par}
    \vspace{0.8cm}
    {\usebeamercolor[fg]{institute}\normalsize\insertinstitute\par}
    \vspace{0.2cm}
    {\usebeamercolor[fg]{date}\normalsize\insertdate\par}
    \vfill
  \end{centering}
}

% Custom command to show navigation symbols after title and outline pages
\makeatletter
\def\shownavigationsymbols{%
    \ifnum\c@framenumber>2%
        \setbeamertemplate{navigation symbols}[default]%
    \else%
        \setbeamertemplate{navigation symbols}{}%
    \fi%
}
\addtobeamertemplate{footline}{\shownavigationsymbols}{}
\makeatother

% Customize frame title
\setbeamertemplate{frametitle}{
    \insertframetitle
}

% Customize itemize bullets
\setbeamertemplate{itemize item}{\small\raise0.5pt\hbox{\textbullet}}
\setbeamertemplate{itemize subitem}{\tiny\raise1.5pt\hbox{\textbullet}}

% Packages
\usepackage{graphicx}
\usepackage{amsmath}
\usepackage{hyperref}
\usepackage{multirow}
\usepackage{array}
\usepackage{comment}

% Paper details (replace with your information)
\title{Coffee}
\date{\today}


% Mark the page number (without any other info) in the footer except the title page
\setbeamertemplate{footline}{%
    \ifnum\value{framenumber}>1
    \hfill\insertframenumber\fi
}


\begin{document}

% Title page
\begin{frame}
    \titlepage
\end{frame}

% ============================================================
%  Frame 1: Conceptualizing "Displacement"
% ============================================================
\begin{frame}{Conceptualizing Displacement: Active Purchase Goal}
\footnotesize

\textbf{Core Idea.}
The displacement indicator $D_{ie}$ is a \emph{revealed-preference proxy} for an
\textbf{active purchase goal}---a chronically activated consumption intention
made visible by customers' purchasing pattern.

\vspace{0.5cm}

\medskip
\tiny
\begin{tabular}{@{} p{0.30\linewidth} p{0.64\linewidth} @{}}
\toprule
\textbf{Concept} & \textbf{Key Statement \& Citation} \\
\midrule
Habitual goal pursuit &
Habitual behavior is fundamentally goal-driven; recurring
purchase routines reflect a chronically active consumption goal,
not mere inertia.
\newline
\textit{Kruglanski \& Szumowska (2020, Perspectives on Psychological Science)} \\
\addlinespace[4pt]
 &
Goals are mental representations that, once activated,
remain accessible and direct behavior through a network
of means--ends associations.
\newline
\textit{Kruglanski et al.\ (2002, Advances in Experimental Social Psychology)} \\
\addlinespace[4pt]
Implemental mindset &
An initial purchase shifts the consumer from a deliberative
to an implemental mindset, making subsequent purchases more
likely---consecutively purchasing customers are already in
this action-oriented state.
\newline
\textit{Dhar, Huber \& Khan (2007, Journal of Marketing Research);
Gollwitzer (1990)} \\
\bottomrule
\end{tabular}
\end{frame}


% ============================================================
%  Frame 2: Why Closure Affects Displaced Customers
% ============================================================
\begin{frame}{Conceptualize Closure Impact on Displaced Customers}
\footnotesize

\textbf{Core Idea.}
When a store closure blocks an active purchase goal, the goal
persists cognitively and intensifies, producing a narrowly-targeted
rebound upon reopening.

\vspace{0.5cm}

\medskip
\tiny
\begin{tabular}{@{} p{0.26\linewidth} p{0.68\linewidth} @{}}
\toprule
\textbf{Mechanism} & \textbf{Key Statement \& Citation} \\
\midrule
Goal persistence (Zeigarnik effect) &
Unfulfilled goals remain mentally active---causing intrusive
thoughts and heightened accessibility---until they are completed
or a concrete plan is formed.
\newline
\textit{Masicampo \& Baumeister (2011, JPSP)} \\
\addlinespace[4pt]
Mental simulation \& counterfactual thinking &
Blocked goals trigger counterfactual and prefactual thinking
(``I would be buying\ldots''), which is closely connected to
goal cognition and regulates subsequent behavior via both
content-specific and motivational pathways.
\newline
\textit{Epstude \& Roese (2008, PSPR);
Kappes (2016, SPPC)} \\
\addlinespace[4pt]
Psychological reactance &
Restricting consumers' freedom to choose triggers reactance---a
motivational state that intensifies desire for the restricted
option and drives efforts to restore the threatened freedom.
\newline
\textit{Brehm (1966); Clee \& Wicklund (1980, JCR)} \\
\addlinespace[4pt]
Consumer response to unavailability &
As personal commitment to an unavailable option increases,
consumers react more negatively and are more motivated
to fulfill the thwarted purchase.
\newline
\textit{Fitzsimons (2000, JCR)} \\
\bottomrule
\end{tabular}
\end{frame}

\end{document}
